
\begin{lemma}
\label{lemma-nested-quotients}
Let $\mathcal{D}$ be a triangulated category and let
$\mathcal{B} \subset \mathcal{A} \subset \mathcal{D}$ be full
triangulated subcategories.
\begin{enumerate}
\item The multiplicative system used to form
$\mathcal{A}/\mathcal{B}$ is the intersection with
$\text{Arrows}(\mathcal{A})$ of the multiplicative system used to form
$\mathcal{D}/\mathcal{B}$.
\item The canonical exact functor
$$
\mathcal{A}/\mathcal{B} \longrightarrow \mathcal{D}/\mathcal{B}
$$
is fully faithful. It is injective on objects, and its image is the full
triangulated subcategory on the objects coming from $\mathcal{A}$.
\item After identifying $\mathcal{A}/\mathcal{B}$ with this image, the
canonical exact functor
$$
\mathcal{D}/\mathcal{A} \longrightarrow
(\mathcal{D}/\mathcal{B})/(\mathcal{A}/\mathcal{B})
$$
is an equivalence of triangulated categories.
\end{enumerate}
\end{lemma}

\begin{proof}
Write $S_{\mathcal{D}}(\mathcal{B})$ and
$S_{\mathcal{A}}(\mathcal{B})$ for the two multiplicative systems in
question. By Lemma \ref{lemma-construct-multiplicative-system}, an arrow
belongs to either system precisely when it occurs as the first arrow of a
distinguished triangle whose third object is isomorphic to an object of
$\mathcal{B}$. The distinguished triangles of $\mathcal{A}$ are exactly
the distinguished triangles of $\mathcal{D}$ all of whose objects belong
to $\mathcal{A}$; see Lemma \ref{lemma-triangulated-subcategory}. This
proves (1).

Let $X$ be an object of $\mathcal{A}$. We claim that the full subcategory
of $S_{\mathcal{D}}(\mathcal{B})/X$ consisting of arrows whose source is
in $\mathcal{A}$ is cofinal. Indeed, suppose that
$s : X' \to X$ belongs to $S_{\mathcal{D}}(\mathcal{B})$. Complete $s$
to a distinguished triangle
$$
X' \longrightarrow X \longrightarrow B \longrightarrow X'[1]
$$
where $B$ is isomorphic to an object of $\mathcal{B}$. Since
$\mathcal{B} \subset \mathcal{A}$ and $\mathcal{A}$ is a full
triangulated subcategory, $X'$ is isomorphic to an object $X''$ of
$\mathcal{A}$. After precomposing $s$ by such an isomorphism we obtain an
object of $S_{\mathcal{A}}(\mathcal{B})/X$ isomorphic to $s$ in
$S_{\mathcal{D}}(\mathcal{B})/X$. The claim follows.

For $X,Y$ in $\mathcal{A}$, Categories, Remark
\ref{categories-remark-right-localization-morphisms-colimit}, fullness of
$\mathcal{A}$ in $\mathcal{D}$, and the cofinality just proved give
\begin{align*}
\Mor_{\mathcal{D}/\mathcal{B}}(X,Y)
&=
\mathop{\rm colim}_{s:X'\to X\text{ in }S_{\mathcal{D}}(\mathcal{B})}
\Mor_{\mathcal{D}}(X',Y) \\
&=
\mathop{\rm colim}_{s:X'\to X\text{ in }S_{\mathcal{A}}(\mathcal{B})}
\Mor_{\mathcal{A}}(X',Y) \\
&=
\Mor_{\mathcal{A}/\mathcal{B}}(X,Y).
\end{align*}
This proves (2). Exactness follows from
Lemma \ref{lemma-universal-property-quotient}.

The composite
$$
\mathcal{D} \longrightarrow \mathcal{D}/\mathcal{B}
\longrightarrow (\mathcal{D}/\mathcal{B})/(\mathcal{A}/\mathcal{B})
$$
annihilates $\mathcal{A}$, and hence factors through an exact functor from
$\mathcal{D}/\mathcal{A}$. Conversely, the quotient functor
$\mathcal{D} \to \mathcal{D}/\mathcal{A}$ annihilates $\mathcal{B}$,
so it first factors through $\mathcal{D}/\mathcal{B}$; this factorization
annihilates $\mathcal{A}/\mathcal{B}$ and therefore factors through
$(\mathcal{D}/\mathcal{B})/(\mathcal{A}/\mathcal{B})$. The uniqueness in
Lemma \ref{lemma-universal-property-quotient} shows that the resulting
two functors are quasi-inverse. This proves (3).
\end{proof}

\begin{lemma}
\label{lemma-subcategories-of-quotient}
Let $\mathcal{D}$ be a triangulated category, let
$\mathcal{B} \subset \mathcal{D}$ be a full triangulated subcategory,
and let
$$
Q : \mathcal{D} \longrightarrow \mathcal{D}/\mathcal{B}
$$
be the quotient functor. Set $\mathcal{K}=\Ker(Q)$.
\begin{enumerate}
\item If $\mathcal{C}$ is a full triangulated subcategory of
$\mathcal{D}/\mathcal{B}$, then
$$
Q^{-1}(\mathcal{C})=
\{X\in\Ob(\mathcal{D})\mid Q(X)\text{ is isomorphic to an object of }
\mathcal{C}\}
$$
is a strictly full triangulated subcategory of $\mathcal{D}$. Moreover,
the multiplicative system associated to $Q^{-1}(\mathcal{C})$ is
$$
\{f\in\text{Arrows}(\mathcal{D})\mid
Q(f)\text{ belongs to the multiplicative system associated to }
\mathcal{C}\}.
$$
\item The rules
$$
\mathcal{A}\longmapsto \overline{Q(\mathcal{A})}
\quad\text{and}\quad
\mathcal{C}\longmapsto Q^{-1}(\mathcal{C})
$$
are mutually inverse, inclusion-preserving bijections between
\begin{enumerate}
\item strictly full triangulated subcategories $\mathcal{A}$ of
$\mathcal{D}$ which contain $\mathcal{K}$, and
\item strictly full triangulated subcategories $\mathcal{C}$ of
$\mathcal{D}/\mathcal{B}$.
\end{enumerate}
Here $\overline{Q(\mathcal{A})}$ denotes the strictly full closure of the
image of $\mathcal{A}$.
\item For every full triangulated subcategory $\mathcal{C}$ of
$\mathcal{D}/\mathcal{B}$, the subcategory $Q^{-1}(\mathcal{C})$ is
saturated if and only if $\mathcal{C}$ is saturated. Equivalently, under
the bijection in (2), $\mathcal{A}$ is saturated if and only if
$\overline{Q(\mathcal{A})}$ is saturated. In particular, if
$\mathcal{A}$ is a full saturated triangulated subcategory of
$\mathcal{D}$ containing $\mathcal{B}$, then its image in
$\mathcal{D}/\mathcal{B}$ is saturated.
\end{enumerate}
\end{lemma}

\begin{proof}
Since $Q$ is exact, inverse images under $Q$ preserve shifts and
distinguished triangles. Strict fullness follows from the definition of
$Q^{-1}(\mathcal{C})$. This proves the first assertion of (1). If
$(X,Y,Z,f,g,h)$ is a distinguished triangle of $\mathcal{D}$, then its
image under $Q$ is distinguished. Thus the cone of $f$ is isomorphic to
an object of $Q^{-1}(\mathcal{C})$ if and only if the cone of $Q(f)$ is
isomorphic to an object of $\mathcal{C}$. By
Lemma \ref{lemma-construct-multiplicative-system}, this proves the final
assertion of (1).

Every inverse image in (2) contains $\mathcal{K}$. Conversely, if
$\mathcal{A}$ is strictly full, triangulated, and contains
$\mathcal{K}$, then $\mathcal{B}\subset\mathcal{A}$ and
Lemma \ref{lemma-nested-quotients} shows that its image has a strictly
full triangulated closure in $\mathcal{D}/\mathcal{B}$. It is immediate
that
$$
\overline{Q(Q^{-1}(\mathcal{C}))}=\mathcal{C}.
$$
We prove the other equality. Suppose that $X$ is an object of
$Q^{-1}(\overline{Q(\mathcal{A})})$. There is an object $Y$ of
$\mathcal{A}$ and an isomorphism $Q(Y)\to Q(X)$. Represent this
isomorphism by a fraction
$$
Q(Y) \xleftarrow{Q(s)} Q(Z) \xrightarrow{Q(f)} Q(X),
$$
where $Q(s)$ and $Q(f)$ are isomorphisms. Categories, Lemma
\ref{categories-lemma-what-gets-inverted} and
Lemma \ref{lemma-operations} show that $s$ and $f$ belong to the
multiplicative system associated to $\mathcal{K}$. Hence they may be
completed to distinguished triangles whose third objects are isomorphic
to objects of $\mathcal{K}$. Since $Y$ and $\mathcal{K}$ are contained
in $\mathcal{A}$, the first triangle gives $Z\in\Ob(\mathcal{A})$; the
second then gives $X\in\Ob(\mathcal{A})$. Here we use strict fullness of
$\mathcal{A}$ to replace objects by isomorphic ones. Therefore
$$
Q^{-1}(\overline{Q(\mathcal{A})})=\mathcal{A},
$$
which proves (2).

Replacing $\mathcal{C}$ by its strictly full closure changes neither
$Q^{-1}(\mathcal{C})$ nor whether $\mathcal{C}$ is saturated. Thus, for
the first assertion of (3), we may assume that $\mathcal{C}$ is strictly
full and put $\mathcal{A}=Q^{-1}(\mathcal{C})$. Part (2) gives
$\mathcal{C}=\overline{Q(\mathcal{A})}$. Suppose first that
$\mathcal{A}$ is saturated and that
$U\oplus V$ is isomorphic to an object of
$\overline{Q(\mathcal{A})}$. Since the quotient functor is the identity
on objects, write $U=Q(X)$ and $V=Q(Y)$. Part (2) gives
$X\oplus Y\in\Ob(\mathcal{A})$. Saturation and strict fullness of
$\mathcal{A}$ imply that $X,Y\in\Ob(\mathcal{A})$, and hence $U,V$ lie
in $\overline{Q(\mathcal{A})}$. Thus the latter is saturated. The converse
follows similarly: if $X\oplus Y$ belongs to $Q^{-1}(\mathcal{C})$ and
$\mathcal{C}$ is saturated, then $Q(X)$ and $Q(Y)$ are isomorphic to
objects of
$\mathcal{C}$, so $X$ and $Y$ belong to $Q^{-1}(\mathcal{C})$. This
proves the equivalence in (3).

Finally, let $\mathcal{A}$ be full, saturated, and contain $\mathcal{B}$.
Its strictly full closure $\overline{\mathcal{A}}$ is saturated and
contains $\mathcal{K}$ by Lemma \ref{lemma-operations}. The case just
proved shows that $\overline{Q(\overline{\mathcal{A}})}$ is saturated.
This is the strictly full closure of the image of $\mathcal{A}$, and the
final assertion follows.
\end{proof}
